\section{Conclusion}
\label{sec:conclusion}

We introduced \textbf{AnomalyClaw}, a training-free visual anomaly detection agent whose headline design is an \emph{always-on parallel-Direct branch}: a single per-item invocation runs a refutation-style tool-using trajectory and a descriptor-free Direct VLM call simultaneously on the same backbone, then returns their unweighted $0.5{-}0.5$ average. On \textbf{CrossDomainVAD-12} ($1{,}418$ test items across $12$ domains) the ensemble is positive-in-mean on all three backbones and statistically significant at $95\%$ on two of three: GPT-5.4 macro AUROC $0.731\!\to\!0.772$ ($+4.01$ pp, $P{=}1.000$), SeedVL $0.678\!\to\!0.738$ ($+5.99$ pp, $P{=}1.000$), Qwen3.5-VL-27B $0.714\!\to\!0.732$ ($+1.69$ pp, $P{=}0.971$; CI $[-0.08,+3.51]$ touches zero). The gain is produced by complementary errors between Direct and the agent, which persist even when the agent is individually weaker than Direct: on Qwen3.5 the agent alone is significantly worse than Direct ($-4.46$ pp, $P{=}0.001$), and yet the ensemble still improves over Direct. This \emph{failure-mode robustness} is the load-bearing property of the design.

Supporting findings from an earlier router/fusion design exploration (Findings~1--4, evaluated on an equivalent twelve-domain split of the same source datasets) are retained for context: task-anchored domain descriptors dominate any agentic intervention at zero inference cost; a calibration-argmax router beats every single-strategy baseline on every backbone; and six failed-variant ablations show that \emph{unconditional} second-call escalation hurts when the baseline is already strong. The parallel-Direct branch side-steps the sequential-rationalisation failure mode of those ablations by making Direct and agent independent score sources rather than successive reasoning turns.

\paragraph{Limitations.}
Three domains remain difficult across all three backbones: D6 Real3D-AD (rendered 3D point-cloud views, $0.47$--$0.56$ AUROC; genuine modality mismatch with a 2D VLM), D10 BMAD liver CT ($0.42$--$0.49$; item-level sampling difficulty in the test split), and D4 SDNET concrete crack detection on Direct alone ($0.44$--$0.61$; the agent partially recovers). The refutation agent alone is significantly weaker than Direct on Qwen3.5; the ensemble papers over this via the parallel-Direct branch, but a deeper fix would require tool-usage refinement for weaker backbones, not just an ensemble. The ensemble weight is fixed a priori at $\alpha{=}0.5$; a per-backbone calibration-split selector is an obvious extension we leave to future work.

\paragraph{Semi-supervised controller extension.}
When a small dev-label budget is available, \S\ref{sec:controller} adds a Controller VLM on top of the two branches that reads both branches' score--rationale pairs and a retrieved rulebook. On Qwen3.5-VL-27B --- the backbone where the passive ensemble gain just touches $95\%$ significance --- the controller reaches $0.7539$ vs the passive ensemble's $0.7333$ ($+2.05$ pp, $P{=}0.995$), and a branch-frozen four-way ablation with a wrong-domain negative control shows that rule content, not the Controller mechanism or the presence of any text, carries the gain. The controller is a single-backbone case study; the passive ensemble remains our training-free main system and the controller a semi-supervised extension at a $480$-label budget.

\paragraph{Future work.}
(i) Calibration-split $\alpha$ selection: we froze $\alpha{=}0.5$ a priori; a 20-item calibration split per backbone should tighten the SeedVL and GPT gains further. (ii) Tool-usage refinement for weak VLM backbones: the Qwen3.5 agent's $-4.46$~pp deficit vs Direct localises to specific tools (e.g.\ LEVIR \texttt{tool\_side\_by\_side}); a small per-backbone tool-catalogue prune may close the deficit and push the Qwen3.5 ensemble into clear significance. (iii) Extending the parallel-Direct idea to ensembles of \emph{two independent agents} (e.g.\ refutation + descriptor-routed) inside a single per-item call, amortising the two-agent cost by running them concurrently on the same VLM API session. (iv) Transfer of the Controller extension to GPT-5.4 and SeedVL backbones, and (v) online rule evolution as test-time feedback accumulates rather than fixing the rulebook at training time.
